\section{Collecting local operations by party}
\label{sec:peps_party_factorization}

\begin{definition}[The registers of a party]
    \label{def:peps_party_registers}
    \lean{TNLean.PEPS.PairEffect.Word.parties}
    \lean{TNLean.PEPS.PairEffect.Word.owners_mem_parties}
    \lean{TNLean.PEPS.PairEffect.Layout.atParty}
    \lean{TNLean.PEPS.PairEffect.partyLayout}
    \lean{TNLean.PEPS.PairEffect.tensorPartyMaps}
    \lean{TNLean.PEPS.PairEffect.tensorPartyMaps_cons_tmul}
    \lean{TNLean.PEPS.PairEffect.partyLayout_eq_of_atParty_eq}
    \lean{TNLean.PEPS.PairEffect.tensorPartyMaps_heq}
    \lean{TNLean.PEPS.PairEffect.Layout.conj_memCongr_heq}
    \lean{TNLean.PEPS.PairEffect.Layout.norm_conj_memCongr}
    \lean{TNLean.PEPS.PairEffect.Layout.eq_conj_memCongr_of_heq}
    \leanok
    \uses{def:peps_party_layout}
    For an ordered list of registers $\ell$ and a party $p$, let $\ell_p$
    retain precisely the registers owned by $p$, in their original order.
    Write $\mathcal H_{\ell,p}$ for their joint memory. The memory of an
    empty list is $\mathbb C$. A finite ordered list of distinct parties
    $p_1,\ldots,p_s$ gives the grouped memory
    $\bigotimes_{j=1}^s\mathcal H_{\ell,p_j}$. Local maps between these
    memories act by their ordered tensor product; for an empty party list
    this is the identity on $\mathbb C$. Equal register lists are identified
    by the identity on their memories, preserving operator norms.

    The parties of a composition include every register owner and the party
    assigned to each local operation. In particular, a scalar operation
    $\mathbb C\to\mathbb C$ retains its assigned party when both its input
    and output register lists are empty. These are the local spaces and
    operations in \cite[Theorem~5.2]{OpenAI2026PolynomialPEPS},
    \texttt{04-compression.tex}, lines 233--251.
\end{definition}

\begin{definition}[Canonical partitions and grouping]
    \label{def:peps_party_partition}
    \lean{TNLean.PEPS.PairEffect.Layout.restrict}
    \lean{TNLean.PEPS.PairEffect.Layout.restrict_nil}
    \lean{TNLean.PEPS.PairEffect.Layout.restrict_cons}
    \lean{TNLean.PEPS.PairEffect.Layout.restrict_append}
    \lean{TNLean.PEPS.PairEffect.Layout.restrict_eq_self_of_owner}
    \lean{TNLean.PEPS.PairEffect.Layout.restrict_eq_nil_of_owner}
    \lean{TNLean.PEPS.PairEffect.Layout.memCongr}
    \lean{TNLean.PEPS.PairEffect.Layout.memCongr_rfl}
    \lean{TNLean.PEPS.PairEffect.Layout.partitionIso}
    \lean{TNLean.PEPS.PairEffect.Layout.partitionIso_cons_true}
    \lean{TNLean.PEPS.PairEffect.Layout.partitionIso_cons_false}
    \lean{TNLean.PEPS.PairEffect.Layout.partitionIso_cons_true_tmul}
    \lean{TNLean.PEPS.PairEffect.Layout.partitionIso_cons_false_tmul}
    \lean{TNLean.PEPS.PairEffect.Layout.partitionIso_append_true_tmul}
    \lean{TNLean.PEPS.PairEffect.Layout.partitionIso_append_false_tmul}
    \lean{TNLean.PEPS.PairEffect.Layout.partitionIso_cons_cons_true_true_tmul}
    \lean{TNLean.PEPS.PairEffect.Layout.partitionIso_cons_cons_true_false_tmul}
    \lean{TNLean.PEPS.PairEffect.Layout.partitionIso_cons_cons_false_true_tmul}
    \lean{TNLean.PEPS.PairEffect.Layout.partitionIso_cons_cons_false_false_tmul}
    \lean{TNLean.PEPS.PairEffect.Layout.withoutParty}
    \lean{TNLean.PEPS.PairEffect.Layout.atParty_withoutParty}
    \lean{TNLean.PEPS.PairEffect.Layout.mem_withoutParty}
    \lean{TNLean.PEPS.PairEffect.Layout.owners_withoutParty}
    \lean{TNLean.PEPS.PairEffect.partyLayout_withoutParty}
    \lean{TNLean.PEPS.PairEffect.groupByPartyIso}
    \lean{TNLean.PEPS.PairEffect.groupByPartyIso_cons}
    \leanok
    \uses{def:peps_party_registers}
    For a choice of parties $A$, there is a prescribed isometric
    identification
    \begin{align}
        U_{A,\ell}:\mathcal H_\ell
         & \longrightarrow \mathcal H_{\ell|_A}\otimes
        \mathcal H_{\ell|_{A^c}},
        \label{eq:peps_party_partition}
    \end{align}
    which retains the order within each selected list. It is obtained
    solely from associativity, exchanges of tensor factors, and the unit
    identifications with $\mathbb C$.

    If $p_1,\ldots,p_s$ are distinct and contain every owner in $\ell$, successive
    partitions give a prescribed isometric identification
    $U_\ell:\mathcal H_\ell\to
        \bigotimes_{j=1}^s\mathcal H_{\ell,p_j}$.
\end{definition}

\begin{proof}\leanok
    Induct on $\ell$. For the empty list, $U_{A,\emptyset}(z)=1\otimes z$
    on $\mathbb C$. Let $r$ be a register with vector $x$, let
    $y\in\mathcal H_\ell$ with $U_{A,\ell}(y)=a\otimes b$, where
    $a\in\mathcal H_{\ell|_A}$ and $b\in\mathcal H_{\ell|_{A^c}}$.
    For the list $(r,\ell)$ consisting of $r$ followed by $\ell$, define
    $U_{A,(r,\ell)}$ by linearity from
    \begin{align*}
        U_{A,(r,\ell)}(x\otimes y)
        =\begin{cases}
             (x\otimes a)\otimes b & \text{if the owner of $r$ lies in $A$}, \\
             a\otimes(x\otimes b)  & \text{otherwise}.
         \end{cases}
    \end{align*}
    The first case is the inverse associativity map and the second the
    exchange of $x$ past $a$, each applied after
    $\operatorname{id}_r\otimes U_{A,\ell}$. For the second assertion,
    successively separate the registers of $p_1,\ldots,p_s$.
    Distinctness prevents any register from being repeated, and coverage
    makes the final remaining register list empty.
\end{proof}

\begin{theorem}[Restriction to selected parties]
    \label{thm:peps_party_restriction}
    \lean{TNLean.PEPS.PairEffect.Word.castLayouts}
    \lean{TNLean.PEPS.PairEffect.Word.sources_castLayouts}
    \lean{TNLean.PEPS.PairEffect.Word.isAllowed_castLayouts}
    \lean{TNLean.PEPS.PairEffect.Word.eval_castLayouts_heq}
    \lean{TNLean.PEPS.PairEffect.Word.eval_castLayouts}
    \lean{TNLean.PEPS.PairEffect.Word.restrict}
    \lean{TNLean.PEPS.PairEffect.Word.sources_restrict}
    \lean{TNLean.PEPS.PairEffect.Word.isAllowed_restrict}
    \lean{TNLean.PEPS.PairEffect.Word.restrict_localMap_true_eval_heq}
    \lean{TNLean.PEPS.PairEffect.Word.restrict_localMap_false_eval_heq}
    \lean{TNLean.PEPS.PairEffect.Word.restrict_swap_true_true_eval_heq}
    \lean{TNLean.PEPS.PairEffect.Word.restrict_swap_true_false_eval_heq}
    \lean{TNLean.PEPS.PairEffect.Word.restrict_swap_false_true_eval_heq}
    \lean{TNLean.PEPS.PairEffect.Word.restrict_swap_false_false_eval_heq}
    \lean{TNLean.PEPS.PairEffect.Word.restrict_frame_true_eval_heq}
    \lean{TNLean.PEPS.PairEffect.Word.restrict_frame_false_eval_heq}
    \lean{TNLean.PEPS.PairEffect.Word.parties_castLayouts}
    \lean{TNLean.PEPS.PairEffect.Word.parties_restrict}
    \leanok
    \uses{def:peps_party_registers,def:peps_party_partition}
    Let $F$ be a composition containing no pair-source preparations.
    For any choice of parties $A$, retaining their registers and local
    operations gives a composition $F_A$ containing no pair-source
    preparations. It is allowed whenever $F$ is allowed. Its local
    operations occur in their original chronological order, and its parties
    are precisely those of $F$ that belong to $A$.
\end{theorem}

\begin{proof}\leanok
    Retain a local operation exactly when its assigned party belongs to $A$.
    Restrict every exchange to the retained registers, and restrict
    compositions and untouched registers recursively. The local
    contraction bounds are unchanged. Induction on the composition proves
    the assertion about its parties, including scalar local operations.
\end{proof}

\begin{theorem}[Factorization across a partition of the parties]
    \label{thm:peps_party_binary_factorization}
    \lean{TNLean.PEPS.PairEffect.Word.partitionIso_eval}
    \leanok
    \uses{def:peps_party_partition,thm:peps_party_restriction}
    If $F:\mathcal H_\ell\to\mathcal H_{\ell'}$ contains no pair-source
    preparations, then for any choice of parties $A$ its restrictions satisfy
    \begin{align}
        U_{A,\ell'}F & =(F_A\otimes F_{A^c})U_{A,\ell}.
        \label{eq:peps_party_binary_factorization}
    \end{align}
    This identity does not require bounds on the norms of the local maps.
\end{theorem}

\begin{proof}\leanok
    Induct on the composition. Let a local operation $L$ at party $p$ act
    on registers $a$, all owned by $p$, in front of a list $\ell$, and let
    $J_{a,\ell}$ identify the memory of $a$ followed by $\ell$ with
    $\mathcal H_a\otimes\mathcal H_\ell$. For $s\in\mathcal H_a$,
    $x\in\mathcal H_{\ell|_A}$, and $y\in\mathcal H_{\ell|_{A^c}}$,
    \begin{align*}
        U_{A,(a,\ell)}J_{a,\ell}^{-1}
        \bigl(s\otimes U_{A,\ell}^{-1}(x\otimes y)\bigr)
        =\begin{cases}
             J_{a,\ell|_A}^{-1}(s\otimes x)\otimes y
              & \text{if $p\in A$},    \\
             x\otimes J_{a,\ell|_{A^c}}^{-1}(s\otimes y)
              & \text{if $p\notin A$}.
         \end{cases}
    \end{align*}
    Thus $L$ acts on $s$ inside the first factor when $p\in A$ and inside
    the second otherwise, while the other restriction is the identity.
    An exchange of two registers is checked in the same way in the four
    cases determined by which of the two registers have owners in $A$; it
    either stays within one part or changes only the prescribed ordering
    between the parts. Tensoring with an untouched register and
    composition preserve the identity.
\end{proof}

\begin{theorem}[Tensor product of the local maps]
    \label{thm:peps_party_factorization}
    \lean{TNLean.PEPS.PairEffect.norm_tensorPartyMaps_le_one}
    \lean{TNLean.PEPS.PairEffect.Word.exists_tensorPartyMaps}
    \lean{TNLean.PEPS.PairEffect.Word.exists_tensorPartyMaps_parties}
    \lean{TNLean.PEPS.PairEffect.Word.eq_nil_and_eval_eq_id_of_parties_eq_empty}
    \leanok
    \uses{def:peps_party_registers,def:peps_party_partition}
    Let $F:\mathcal H_\ell\to\mathcal H_{\ell'}$ be an allowed composition
    containing no pair-source preparations. Let $p_1,\ldots,p_s$ be distinct
    parties containing every party of this composition. There are
    contractions $B_{p_j}:\mathcal H_{\ell,p_j}\to
        \mathcal H_{\ell',p_j}$ such that
    \begin{align}
        U_{\ell'}F
         & =\left(\bigotimes_{j=1}^s B_{p_j}\right)U_\ell.
        \label{eq:peps_party_factorization}
    \end{align}
    The tensor product is itself a contraction. The list of parties may be
    chosen to consist exactly of the parties of $F$. Each map is obtained
    by composing the local operations at its party. A composition without
    parties has empty input and output register lists and is the identity
    on $\mathbb C$; this is the case $s=0$.
    This is the collection of local maps in
    \cite[proof of Theorem~5.2]{OpenAI2026PolynomialPEPS},
    \texttt{04-compression.tex}, lines 233--251.
\end{theorem}

\begin{proof}\leanok
    \uses{thm:peps_party_binary_factorization,thm:peps_party_restriction}
    Apply (\ref{eq:peps_party_binary_factorization}) successively with
    $A=\{p_j\}$. The remaining
    composition has no parties once all $p_j$ have been removed and is the
    identity on $\mathbb C$. The restriction to each party is allowed, so
    its operator has norm at most one. The tensor product of the resulting
    contractions gives (\ref{eq:peps_party_factorization}).
\end{proof}

\begin{theorem}[Prepared pair sources and local contractions]
    \label{thm:peps_prepared_party_maps}
    \lean{TNLean.PEPS.PairEffect.Word.exists_prepared_tensorPartyMaps}
    \leanok
    \uses{def:peps_source_inventory,def:peps_party_registers,def:peps_party_partition}
    Let $P$ be the prescribed finite set of parties of a gate, and use a
    fixed ordering of $P$. Let $F:\mathcal H_{\rm in}\to\mathcal H_{\rm out}$ be
    an allowed composition on registers owned by these parties. There is a
    normalized source list $S$, containing exactly one source for each
    unordered source pair occurring in $F$, and a contraction $B_p$ at each $p\in P$
    such that
    \begin{align}
        U_{\rm out}F
         & =\left(\bigotimes_{p\in P}B_p\right)
        U_{S,\rm in}P_{S,\mathcal H_{\rm in}}.
        \label{eq:peps_prepared_party_maps}
    \end{align}
    Here $U_{S,\rm in}$ groups the prepared source registers together with
    the original input registers by their owners; $U_{\rm out}$ groups the
    output registers in the same party order. The input and output of $B_p$
    are precisely the corresponding memories of party $p$.
\end{theorem}

\begin{proof}\leanok
    \uses{thm:peps_grouped_sources,thm:peps_party_factorization}
    Theorem~\ref{thm:peps_grouped_sources} gives
    $F=V\circ P_{S,\mathcal H_{\rm in}}$ with $V$ allowed and containing no
    sources. Every party named by $V$ belongs to the prescribed set $P$.
    Apply (\ref{eq:peps_party_factorization}) to $V$ with all parties of $P$,
    and compose on the right with the preparation map.
\end{proof}
